\chapter{Reinforcement learning in ten minutes}
\label{ch:rl}

This chapter gives just enough reinforcement learning (RL) to make every
number in the tricycle configuration meaningful. If you know PPO already, skim
Section~\ref{sec:rl_numbers} for how the general ideas map onto this task.

\section{The loop}

An \textbf{environment} is a simulation that has a state, accepts an
\textbf{action}, advances one \textbf{step}, and returns three things: an
\textbf{observation} (what the agent is allowed to see), a \textbf{reward} (one
number saying how good that step was) and a \textbf{done} flag (whether the
episode is over). An \textbf{agent} is a function, the \textbf{policy}, that
maps observations to actions. Training means adjusting the policy so that the
sum of rewards over an \textbf{episode} (a sequence of steps that ends when
done is true) becomes as large as possible.

For the tricycle: the state is the pose and velocity of the car, the joint
angles and speeds, and the position of every soil particle. The observation is
a 15-number vector (Chapter~\ref{ch:env}). The action is two numbers in
$[-1,1]$: throttle and steering. The reward is how far forward the car moved
during this step, in metres. An episode lasts 2~seconds or until the car
flips, drives off the strip, or the physics blows up.

\section{Vectorised environments}

One simulation step of one car teaches the agent very little. Isaac Lab
therefore simulates $N$ identical copies of the scene at once, called
\textbf{environments} or \textbf{envs}, laid out on a grid in the same physics
world, and returns observations as a tensor of shape $(N, 15)$, rewards as
$(N,)$, and so on. Every method you write in \code{TricycleEnv} operates on
whole tensors; there is never a Python loop over envs. Each env has its own
\textbf{origin} (\code{scene.env_origins}), and positions are usually measured
relative to it so that env 7 does not look different from env 0 just because it
sits further along the grid.

The number of envs is a memory decision. On the 6~GB laptop card the soil
solver limits it to a handful (Chapter~\ref{ch:physics}); on a large GPU the
same code runs thousands.

\section{Policy and value networks}
\label{sec:networks}

The policy is a small neural network. In this task it has two hidden layers of
64 units with the ELU activation. Given the 15 observations it outputs the
\emph{mean} of a Gaussian distribution over the two actions; a second learned
parameter, the standard deviation (initially $1.0$, see
\code{init_noise_std}), controls how much the agent explores around that mean.
During training actions are \emph{sampled} from this Gaussian, which is where
exploration comes from. During play the mean is used.

A second network of the same shape, the \textbf{critic} or \textbf{value
function}, predicts the expected sum of future rewards from a given
observation. It is not used to act; it is used to judge whether an action
turned out better or worse than expected. That difference is the
\textbf{advantage}. Both networks are trained together, which is why the
architecture is called \emph{actor-critic}.

\section{Discounting and GAE}

Rewards far in the future are worth less than rewards now. The \textbf{discount
factor} $\gamma$ (\code{gamma = 0.99}) multiplies each future reward by
$\gamma^k$ after $k$ steps. With $\gamma=0.99$ the effective horizon is about
$1/(1-\gamma) = 100$ steps, which at 50~Hz is exactly one 2-second episode.

The advantage estimate mixes the critic's predictions with actual rewards.
\textbf{Generalised Advantage Estimation} (GAE) blends one-step and many-step
estimates with a second parameter $\lambda$ (\code{lam = 0.95}): higher
$\lambda$ trusts observed rewards more (less bias, more noise), lower trusts
the critic more.

\section{PPO in plain words}

\textbf{Proximal Policy Optimisation} (PPO) is the learning rule used here.
Each \textbf{iteration}:

\begin{enumerate}[leftmargin=1.8em]
  \item \textbf{Roll out.} Run the current policy for
        \code{num_steps_per_env} steps in every env, storing observations,
        actions, rewards, dones and the critic's values. With 4 envs and 50
        steps that is 200 samples per iteration.
  \item \textbf{Compute advantages} with GAE over the stored rollout.
  \item \textbf{Update.} Split the rollout into \code{num_mini_batches}
        mini-batches and pass over them \code{num_learning_epochs} times.
        For each mini-batch, compute the ratio between the probability the
        \emph{new} policy assigns to the stored action and the probability the
        \emph{old} policy assigned. Multiply by the advantage, but
        \emph{clip} the ratio to $[1-\epsilon, 1+\epsilon]$ with
        $\epsilon=$ \code{clip_param} $=0.2$. The clipping is what makes PPO
        ``proximal'': the policy is not allowed to move far from the one that
        collected the data, which keeps learning stable.
  \item Add the critic's loss (mean squared error to the returns, weighted by
        \code{value_loss_coef}) and subtract an \textbf{entropy bonus}
        (\code{entropy_coef}) that rewards keeping the action distribution
        wide, i.e.\ keeps exploring.
  \item Take a gradient step with Adam at \code{learning_rate}, after clipping
        the gradient norm to \code{max_grad_norm}.
\end{enumerate}

With \code{schedule = "adaptive"}, \code{rsl_rl} also measures the KL
divergence between the old and new policies after each mini-batch. If it
exceeds $2\times$\code{desired_kl} the learning rate is divided by 1.5; if it is
below half of \code{desired_kl} the rate is multiplied by 1.5 (bounded to
$[10^{-5}, 10^{-2}]$). This automatic tuning is why the initial
\code{learning_rate = 1e-3} does not need to be exactly right.

\section{Terminated versus timed out}
\label{sec:term_vs_timeout}

An episode can end for two very different reasons and RL algorithms must know
which one happened:

\begin{itemize}[leftmargin=1.6em]
  \item \textbf{Terminated}: the task itself ended (the car flipped over). The
        value of the final state is genuinely zero; nothing good can follow.
  \item \textbf{Timed out} (also called truncated): the clock ran out while the
        car was still driving happily. Its state still has value; the
        episode was merely cut for bookkeeping.
\end{itemize}

Isaac Lab's \code{_get_dones} therefore returns \emph{two} boolean tensors,
\code{terminated} and \code{time_out}, and the \code{rsl_rl} wrapper passes the
time-outs to the learner separately so that it can \emph{bootstrap}: add the
critic's estimate of the remaining value back into the return for timed-out
episodes. Getting this wrong makes the agent believe that surviving to 2
seconds is bad, which is the opposite of what we want.

\section{How the numbers map onto the tricycle}
\label{sec:rl_numbers}

\begin{center}
\small
\begin{tabular}{@{}lll@{}}
\toprule
Quantity & Value & Where it is set \\
\midrule
observation size & 15 & \code{observation_space} in \code{TricycleEnvCfg} \\
action size & 2 (throttle, steer) & \code{action_space} \\
physics step & 10 ms (\code{dt = 1/100}) & \code{SimulationCfg} \\
control step & 20 ms (\code{decimation = 2}) & \code{TricycleEnvCfg} \\
episode length & 2 s $=$ 100 control steps & \code{episode_length_s} \\
reward per step & forward displacement, clipped to $\pm 0.1$ m & \code{_get_rewards} \\
episode return & metres driven forward in 2 s & sum of the above \\
rollout length & 50 steps per env per iteration & \code{num_steps_per_env} \\
iterations & 300 & \code{max_iterations} \\
checkpoint every & 50 iterations & \code{save_interval} \\
networks & $15 \to 64 \to 64 \to 2$ and $15 \to 64 \to 64 \to 1$ & \code{policy} cfg \\
\bottomrule
\end{tabular}
\end{center}

Because the reward is displacement, the ``mean episode reward'' printed by
\code{rsl_rl} is directly readable: a value of 1.4 means the average car drove
1.4~m forward in its 2~s episode. The runbook's success criterion (reward
climbs past $\sim$2~m without diverging) is expressed in the same unit.
